\documentclass[a4paper,12pt]{ThesisStyle}
\usepackage[utf8]{inputenc}
\usepackage{thesis-style}
\usepackage{parskip}

\begin{document}

\frontmatter

\pagenumbering{gobble}

\thispagestyle{empty}
\begin{table}[htb]
\centering
\begin{Large}
\resizebox{\textwidth}{!}{\begin{tabular}{| l |}
  \hline
  \\
\includegraphics[scale=0.9]{imatges/logo_eps.png}\\[0.7cm]
\centerline{Projecte fi de grau}\\[1cm]
\hline
\\
\textbf{Estudi}: Grau en Enginyeria Informàtica\\[0.7cm]
\hline
\\
\textbf{Títol}: Eina de suport per a l’elaboració dels horaris dels graus de l’EPS\\[0.7cm]
\hline
\\
\textbf{Document}: Memòria\\[0.7cm]
\hline
\\
\textbf{Alumne}: Adrià Ribas Chico\\[0.7cm]
\hline
\\
\textbf{Tutor 1}: Dra. Marta Fort Masdevall\\
\textbf{Departament}: Informàtica, matemàtica aplicada i estadística\\
\textbf{Àrea}: Llenguatges i sistemes informàtics\\[0.7cm]
\hline
\\
\textbf{Tutor 2}: Dr. Antonio Rodríguez Benítez\\
\textbf{Departament}: Informàtica, matemàtica aplicada i estadística\\
\textbf{Àrea}: Llenguatges i sistemes informàtics\\[0.7cm]
\hline
\\
\textbf{Convocatòria (mes/any)}: Juny de 2022\\[0.7cm]
\hline

\end{tabular}}
\end{Large}
\end{table}

\newpage

\begin{titlepage}

% Upper part of the page
\includegraphics[scale=0.9]{imatges/logo_eps.png} \\[1cm]
\begin{center}
\textsc{\Large Projecte Fi de Grau} \\[1cm]

% Title
\begin{spacing}{2}
\HRule \\
\textbf{\Huge Eina de suport per a l’elaboració dels horaris dels graus de l’EPS} \\
\HRule \\[0.5cm]
\end{spacing}

% Author and supervisor and other data
{
\large
\emph{Autor:} \\
Adrià \textsc{Ribas Chico} \\[1cm]
Juny 2022 \\[1cm]
Grau en Enginyeria Informàtica \\[1cm]
\emph{Tutors:} \\
Dra. Marta \textsc{Fort Masdevall} \\
Dr. Antonio \textsc{Rodríguez Benítez} \\
}

\end{center}
\end{titlepage}

\titlepage

\dominitoc

\pagenumbering{roman}

\chapter*{Resum}
\label{cap:resum}

El resum del projecte va aquí \ldots

\chapter*{Agraïments}
\label{cap:agraiments}

Per començar vull agrair molt especialment a \ldots


\tableofcontents

%\listoffigures

%\listoftables

\mainmatter

\chapter{Introducció}
\label{cap:intro}

\section{Antecedents}
\label{sec:antecedents}

En aquesta secció, es resumiran els antecedents que han donat peu al plantejament d'aquest projecte.

La idea del projecte neix de determinades necessitats que cert personal de l'Escola Politècnica Superior de la Universitat de Girona fa temps que té. Més concretament, es tracta d'una necessitat del personal encarregat de gestionar tot el que fa referència a la confecció i manteniment dels horaris del centre: horaris dels graus, dels professors, ocupació d'aules i espais, etc.

Actualment, per dur a terme la creació dels horaris, aquestes persones utilitzen mètodes i eines incòmodes i poc àgils, a part de no estar automatitzades ni específicament dissenyades per abordar aquest tipus de tasques.

Ara per ara, per exemple, no tenen manera de detectar incompatibilitats horàries ni solapaments de manera automàtica. A la secció~\ref{sec:situacio_actual} es descriuran els mètodes emprats actualment amb més profunditat.

\section{Situació actual}
\label{sec:situacio_actual}

En aquesta secció, es detallaran quins són i com funcionen els mètodes de gestió i elaboració d'horaris de l'EPS que s'utilitzen actualment. Es veuran quins càrrecs de l'escola hi participen i les seves funcions, així com la manera en què interactuen entre ells i intercanvien la informació corresponent.

Parlar del fluxe dels processos, etc \ldots

Rectorat reparteix crèdits a les facultats. Cada facultat decideix quants grups grans, petits, etc. de cada grau en funció dels alumnes. Introdueix els excels.
A partir d'aquí, entren els coordinadors, per fer els horaris amb els seus mètodes. Canviar-los el mínim possible respecte l'any passat, per intentar evitar
solapaments. Un cop fets, els coordinadors passen els horaris a direcció. S'entren al sistema i revisen solapaments. Si tot està bé, els responsables de docència,
que assignen professors als grups. 

\section{Propòsit}
\label{sec:proposit}

En aquesta secció, s'exposarà el propòsit general del projecte, tenint en compte els antecedents vists a la secció~\ref{sec:antecedents}.

En definitiva, la gestió dels horaris de l'EPS suposa una inversió de temps massa elevada per a la gent que se n'ocupa. És per això que la Dra. Marta Fort Masdevall, coordinadora d'estudi del Grau en Enginyeria Informàtica de la universitat, proposa un projecte de fi de grau que té la finalitat de trobar una solució al problema.

El propòsit és desenvolupar una eina de suport informàtic que permeti a l'usuari elaborar horaris de forma àgil, eficient i segura. L'eina també hauria de ser capaç de comprovar automàticament la disponibilitat de les aules, les incompatibilitats horàries dels professors i la concordança entre les assignatures i el nombre de grups previstos. A més a més, hauria d'oferir diverses vistes per tal que l'usuari pugui visualitzar la informació, com ara:
\begin{itemize}
  \item Els horaris d'un grau per curs i quadrimestre.
  \item Els horaris d'un professor per quadrimestre.
  \item L'ocupació d'un espai per quadrimestre.
\end{itemize}

També es planteja la possibilitat de disposar d'un sistema de control d'usuaris. Cada usuari tindria assignat un conjunt de rols determinat. Els rols representarien els diferents càrrecs del personal de l'EPS en matèria de gestió d'horaris. Així doncs, cada usuari podria executar les accions i consultar la informació que el seu conjunt de rols li permeti. D'aquesta manera, es dividirien les diferents tasques i processos entre rols d'usuari i cadascun dels càrrecs podria realitzar la feina que li correspon. La proposta inicial comprèn els rols d'usuari següents:
\begin{itemize}
  \item \texttt{Administrador}: Introdueix les aules i els grups previstos.
  \item \texttt{Coordinador}: Elabora els horaris.
  \item \texttt{Director de departament}: Dóna d'alta responsables de docència.
  \item \texttt{Responsable de docència}: Dóna d'alta i assigna professors als grups.
  \item \texttt{Professor}: Visualitza el seu horari.
\end{itemize}

A més a més, l'aplicació hauria de ser accessible via web. D'aquesta manera, tots els usuaris podrien utilitzar-la des de qualsevol lloc i dispositiu, sense preocupar-se d'instal·lacions ni actualitzacions.

\section{Motivacions}
\label{sec:motivacions}

En aquesta secció, es parlarà en primera persona sobre quines motivacions personals hi ha darrere del projecte i en justifiquen l'elecció.

\ldots

\section{Objectius generals}
\label{sec:objectius_generals}

En aquesta secció, s'enumeraran els objectius generals del projecte, que ja s'han deixat entreveure prèviament a la secció~\ref{sec:proposit}. No obstant això, a continuació se'n presenta la llista completa:
\begin{itemize}
  \item Proporcionar una interfície còmoda i intuïtiva per dur a terme les tasques de gestió i elaboració dels horaris de l'EPS.
  \item Emmagatzemar i processar dinàmicament les dades i relacions referents als diversos graus, cursos, quadrimestres, assignatures, grups, espais, professors, etc.
  \item Detectar i evitar automàticament qualsevol tipus d'inconsistència o incompatibilitat horària, per tal d'aportar seguretat al treball.
  \item Possibilitar la pujada d'arxius externs de dades que serveixin per obtenir la informació bàsica necessària per al funcionament de l'aplicació i generar possibles punts de partida per a la planificació dels horaris.
  \item Permetre la visualització de l'ocupació de les aules, dels horaris dels professors i dels horaris dels grups de cada grau, entre d'altres vistes que puguin ser d'utilitat pels usuaris.
  \item Admetre diferents rols d'usuari, als quals s'assigni una sèrie de tasques i un conjunt de permisos concret.
\end{itemize}

Al capítol~\ref{cap:requisits} es desenvoluparan aquests objectius generals i es concretaran els requisits específics de l'aplicatiu.


%%%%%%%%%%%%%%%%%%%%%%%%%%%%%%%%%%%%%%%%%%%%%%%%%%%%%%%%%%%%%%%%%%%%%%%%%%%%%%%%%%%%%%%%%%%%%%%%%%%%%%%%%%%%%%%%%%%%%%%%%%%%%%%%%%%%%%%%%%%%%%%%%%%%%
\section{-------------- EXEMPLES I UTILITATS --------------}
\subsection{Paraules per començar seccions}
\textbf{Idees:}

Abordar, concretar, exposar, parlar, descriure, repassar, mostrar, ensenyar, desenvolupar, tractar, veure, aprofundir, investigar, discutir, indagar, detallar,
enumerar,


\subsection{Altres}

Això és un exemple de citació d'un llibre~\cite{Coleman1974}, un article científic~\cite{Ruiz2008} i una referència a una web~\cite{Halcon}.

Exemple de taula:
\begin{table}[htb]
\centering
\begin{tabular}{ | r | c | c | l | }
 \hline
  Any & Matriculats & Aprovats & Percentatge\\
\hline
 2019  & 65 & 47 & 72.3\%\\
 2020  & 69 & 48 & 69.6\%\\
 2021  & 75 & 58 & 77.3\%\\
  \hline
  \end{tabular}
\caption{\label{taula:taulaexemple} Aquí és on s'ha de posar el peu de taula. }
\end{table}

Exemple de figura:
\begin{figure}[htb]
\centering
\includegraphics[width=8 cm]{imatges/logo_eps.png}
\caption{\label{fig:logo} Logotip de l'Escola Politècnica Superior.}
\end{figure}

Exemple de fòrmula:
\begin{equation}
H(X) = -\sum_{i=1}^{N}p_s(x_i) \log \left( p_s(x_i) \right).
\label{equ:entropia}
\end{equation}


També es pot fer referència en el text a les taules (p.ex. veure la Taula~\ref{taula:taulaexemple}), a les figures (p.ex. veure la Figura~\ref{fig:logo}) o a les fòrmules (p.ex. veure Equació~\ref{equ:entropia}).

%%%%%%%%%%%%%%%%%%%%%%%%%%%%%%%%%%%%%%%%%%%%%%%%%%%%%%%%%%%%%%%%%%%%%%%%%%%%%%%%%%%%%%%%%%%%%%%%%%%%%%%%%%%%%%%%%%%%%%%%%%%%%%%%%%%%%%%%%%%%%%%%%


\chapter{Viabilitat}
\label{cap:viabilitat}



\chapter{Metodologia}
\label{cap:metodologia}



\chapter{Marc de treball i conceptes previs}
\label{cap:marcdetreball}
Explicar setmanes a i b, assignatures compartides, què són els blocs (especials també)..., quantes hores
són un crèdit, ¿explicar funcions de cada rol?... 



Explicar tecnologies.

% Grups grans i mitjans no hi ha restriccions d'aules. Grups petits sí,


\chapter{Requisits del sistema}
\label{cap:requisits}

\section{Requisits generals}
\label{sec:requisits_generals} % Poder guardar versions del procés d'elaboració d'horaris.

En aquesta secció, es llistaran els requisits generals del sistema, comuns per a tots els usuaris, independentment dels seus rols.

\begin{itemize}
  \item Assegurar l'autenticació de tots els usuaris a través d'un formulari de \emph{login}. S'ha de mostrar a la pantalla quan un usuari no autenticat accedeix a l'aplicació. Sense estar-ho, no l'ha de poder fer servir.
  \item No permetre el registre d'usuaris des de fora. Usuaris de determinats rols s'encarregaran de donar d'alta altres usuaris del rol que els correspongui.
  \item Possibilitar a tots els usuaris la consulta tant de les seves dades, com les de l'usuari que l'ha donat d'alta, mitjançant perfils personals.
  \item Proporcionar accés als manuals d'usuari de l'aplicació, així com a la informació de contacte i als aspectes legals.
\end{itemize}

% ADMINISTRADORS
\section{Requisits dels Administradors}
\label{sec:requisits_administradors}

En aquesta secció, es llistaran els requisits particulars dels usuaris amb rol d'Administrador.

\begin{itemize}
  \item Carregar el pla docent d'un curs acadèmic per tal d'inicialitzar i poblar la base de dades amb la informació del següent quadrimestre. Aquesta entrada es fa mitjançant un fitxer de format \emph{xlsx} (Microsoft Excel).
    \begin{itemize}
      \item Descriure en algun lloc què porta aquest fitxer i posar-hi una referència des d'aquí?
      \item Què passa si l'Administrador torna a penjar el pla docent quan ja hi hagi horaris fets/començats?
    \end{itemize}
  \item Modificar, si s'escau, un conjunt específic de paràmetres del pla docent sense haver de tornar-lo a carregar, detallat a continuació:
    \begin{itemize}
      \item Nombre d'aules habilitades de cada tipus.
      \item \ldots
    \end{itemize}
  \item Assignar una aula a cada bloc horari, a mesura que els coordinadors acabin d'elaborar els horaris que els pertòquin.
  \item Donar d'alta o de baixa usuaris dels rols Coordinador i Director de departament, a banda de poder veure'n el llistat.
\end{itemize}

% COORDINADORS
\section{Requisits dels Coordinadors}
\label{sec:requisits_coordinadors}

En aquesta secció, es llistaran els requisits particulars dels usuaris amb rol de Coordinador.

\begin{itemize}
  \item Visualitzar els horaris dels graus que coordina.
  \item Visualitzar els horaris dels graus (o cursos o quadrimestres dels graus?) que ofereixin alguna assignatura compartida amb qualsevol dels graus que coordina.
  \item Visualitzar els horaris dels professors que imparteixin docència a alguna de les assignatures que ofereixi qualsevol dels graus que coordina.
  \item Visualitzar els horaris de l'ocupació de qualsevol aula de l'escola.
  \item Carregar opcionalment els horaris que ha tingut anteriorment qualsevol dels graus que coordina per tal d'omplir els actuals i tenir un punt de partida.
  \item Editar els horaris de cada quadrimestre dels cursos de qualsevol dels graus que coordina. El ventall de possibilitats del procés d'edició es detalla a continuació:
    \begin{itemize}
      \item Seleccionar una vista setmanal per tal de visualitzar només els blocs horaris situats a un tipus de setmana concret. Les vistes setmanals són les següents:
        \begin{itemize}
          \item Vista dels blocs horaris situats només a les setmanes de tipus A.
          \item Vista dels blocs horaris situats només a les setmanes de tipus B.
          \item Vista dels blocs horaris situats alhora a tots els tipus de setmana.
          \item Vista de tots els blocs horaris, independentment del tipus de setmana en què estigui situat.
        \end{itemize}
      \item Consultar el nombre d'aules disponibles de qualsevol tipus en una franja horària de setmanes de tipus A i/o B determinada (s'ha d'escriure millor).
      \item Visualitzar els blocs horaris pendents de situar i, si s'escau, modificar-ne la duració.
      \item Situar o moure blocs horaris a franges horàries concretes de setmanes de tipus A i/o B.
      \item Treure blocs horaris ja situats en una franja horària.
      \item Modificar el nombre de blocs horaris de cada tipus de grup de qualsevol assignatura.
      \item Crear nous grups de qualsevol assignatura. Cal que n'especifiqui el tipus i el nombre de blocs horaris. A partir d'això, es crearan els corresponents blocs horaris pendents de situar.
      \item Modificar el número de qualsevol dels grups de les assignatures.
      \item Esborrar grups de qualsevol assignatura.
      \item Crear manualment blocs horaris genèrics que no estiguin associats a cap grup de cap assignatura. Cal que n'especifiqui una duració, un títol i, opcionalment, un subtítol.
    \end{itemize}
\end{itemize}

% DIRECTORS DE DEPARTAMENT
\section{Requisits dels Directors de departament} % Els departaments estan dividits en àrees. Cada àrea té un responsable de docència.
\label{sec:requisits_director_departament}

En aquesta secció, es llistaran els requisits particulars dels usuaris amb rol de Director de departament.

\begin{itemize}
  \item Visualitzar els horaris dels professors que imparteixin docència a alguna assignatura de qualsevol de les àrees del seu departament i dels graus (o curs o quadrimestre del grau?) que en continguin alguna.
  \item Donar d'alta o de baixa usuaris del rol Responsable de docència que s'encarreguin de qualsevol de les àrees del seu departament, a banda de poder veure'n el llistat i consultar-ne el perfil.
\end{itemize}

% RESPONSABLES DE DOCÈNCIA
\section{Requisits dels Responsables de docència} % Està lligat a una àrea. Pot haver-hi més d'un.
\label{sec:requisits_responsables_docencia}

En aquesta secció, es llistaran els requisits particulars dels usuaris amb rol de Responsable de docència.

\begin{itemize} % Pot anar assignant quan els horaris encara s'estan fent
  \item Visualitzar els horaris dels professors que imparteixin docència a alguna assignatura de qualsevol de les seves àrees.
  \item Assignar un professor a cada bloc horari de cadascun dels grups de les assignatures de qualsevol de les seves àrees.
  \item Donar d'alta o de baixa usuaris del rol Professor que hagin d'impartir docència a alguna de les assignatures de la seva àrea, a banda de poder veure'n el llistat i consultar-ne el perfil.
\end{itemize}

% PROFESSORS
\section{Requisits dels Professors}
\label{sec:requisits_professors}

En aquesta secció, es llistaran els requisits particulars dels usuaris amb rol de Professor.

\begin{itemize} % vistes per separat
  \item Visualitzar els seus propis horaris, és a dir, els blocs horaris que li han estat assignats.
  \item Visualitzar els horaris de les assignatures a les que imparteixi docència.
  \item Visualitzar els horaris dels graus (o curs o quadrimestre del grau?) que ofereixin alguna de les assignatures en què imparteixi docència.
\end{itemize}

\chapter{Planificació}  % Ordre canviat. Ordre antic: Planificació --> Marc de treball --> Requisits
\label{cap:planificacio}



\chapter{Estudi i decisions}
\label{cap:estudi}

\begin{itemize}
  \item Visual Studio Code. Extensions?
  \item Git i GitHub
  \item YouTrack ?
  \item Node js
  \item Express
  \item MySQL
  \item Postman, Insomnia o altres?
  \item \ldots
\end{itemize}

\chapter{Anàlisi i disseny del sistema}
\label{cap:analisi}



\chapter{Implementació i proves}
\label{cap:implementacio}



\chapter{Implantació i resultats}
\label{cap:implantacio}



\chapter{Conclusions}
\label{cap:conclusions}



\chapter{Treball futur}
\label{cap:treball_futur}

\begin{itemize}
  \item Poder buscar usuaris per consultar les seves dades (nom, rol, departament, telèfon, email, etc.) i habilitar un link per enviar-li un email.
  \item \ldots
\end{itemize}


\backmatter

\bibliographystyle{ThesisStyleBreakable}
\bibliography{biblio}
%\printnomenclature

%\appendix

%\include{Appendix1}

\chapter*{Manual d'usuari}

\section*{Rol 1}



\section*{Rol 2}



\section*{Rol 3}



\section*{Rol 4}




\end{document}
